% Reúne las reiteraciones del resumen extenso anterior; no sustituye las pruebas.
\subsection*{Carte des résultats et des dépendances}
L’argument comporte cinq niveaux qu’il convient de distinguer à la lecture. La
construction arithmétique fournit les modes ; leurs lecteurs définissent la forme
complète ; les estimations établissent la positivité sur des domaines précis ; le
critère de Weil détermine quelle identification supplémentaire permet de conclure RH.
L’ordre d’annonce rend visible le résultat principal, tandis que le corps
conserve l’ordre de ses dépendances.

\begin{center}
\renewcommand{\arraystretch}{1.22}
\begin{tabular}{@{}p{0.21\linewidth}p{0.48\linewidth}p{0.23\linewidth}@{}}
\hline
\textbf{Niveau} & \textbf{Résultat et fonction dans l'argument} & \textbf{Résidence}\\
\hline
Construction commune &
Histoires, feuillets, mémoire et incidences demeurent dans un même système de raffinement. Centre, frontière relevée et vacances spécifient quelle information transmet chaque lecture. &
Sections~\ref{lectura:manuscrito-05} et~\ref{sec:centro-frontera}.\\
Arithmétique et fonctionnelle &
Produit natif, formes normales et irréductibles ; occupations, produit
d’Euler, moments premier--puissance, parties finies et prolongement thêta--Mellin. &
Sections~\ref{lectura:manuscrito-01}--\ref{lectura:manuscrito-04} ;
début de~\ref{lectura:manuscrito-03}.\\
Reconstruction et transport &
Les colonnes conservent conjointement les termes premier, gamma et polaire.
La queue gamma reconstruit la fonction test. La compression, la mémoire et Cayley
conservent leurs bilans et leurs produits croisés. &
Sections~\ref{lectura:manuscrito-03}, \ref{sec:transporte-momentos} et~\ref{sec:lectores-gamma}.\\
Signe de la forme &
Coercivité de l’assemblage spécifié, contrôle des détails et réduction de
Schur. La fenêtre complète de longueur \(1/9\) vérifie
\(\mathscr W(f,f)\geq\frac12\|f\|_2^2\). &
Section~\ref{lectura:manuscrito-03} et théorèmes~\ref{schur-add-schur-exacto}, \ref{schur-add:coercividad-ventana}.\\
Criterio global &
L'identification des moments complets, avec tous leurs termes croisés et la couverture du domaine de fonctions tests, permettrait d'appliquer le critère de Weil. Elle n'est pas remplacée par les bornes locales précédentes. &
Section~\ref{lectura:manuscrito-03} ;
sous-sections~\ref{tm-add-c5} et~\ref{tm-add-m7}.\\
\hline
\end{tabular}
\end{center}

Deux distinctions accompagnent cette application. En arithmétique, la période d’une action
résiduelle ne remplace pas l’irréductibilité du produit. En analyse, une
factorisation positive des moments de mémoire ne s’identifie pas
automatiquement à celle des moments arithmétiques. Les constantes
\(8/81\) et \(5/9\) appartiennent respectivement au bilan de compression et
au rapport modal de mémoire ; leurs preuves conservent ces fonctions distinctes.

Le centre \((5,5)\), la signature régionale \(555555\), ses trois feuillets de
continuation et les calendriers \(21/22\) et \(6/7\) sont expliqués dans leur position
constructive. Les lectures \(72,27,77,22\), la frontière relevée et la
complétion \(729+271\) ne sont pas regroupées par coïncidence numérique : chacune
conserve l’opération et le domaine qui la produit. Leur fonction dans cet
article est de préciser l’information transportable avant de passer au
problème spectral.

La composition des lecteurs conserve en outre deux opérations distinctes : le produit diagonal des tables résiduelles et le produit tensoriel de leurs moments, qui produit une convolution de diviseurs. Le transport de Stieltjes inclut la contribution affine du pôle. La hiérarchie de Bendersky--Glaisher relie les dérivées régularisées aux niveaux négatifs pairs aux valeurs positives impaires de la même fonctionnelle. Les encadrements par coupure des irréductibles complètent les bornes par profondeur, avec leurs erreurs respectives explicites (sections~\ref{ix-add:corte-primos}, \ref{ix-add:composicion-lectores} et~\ref{ix-add:bendersky}).

Les définitions et les preuves utilisées sont réunies dans les sections indiquées
et dans les appendices A--C, avec les domaines de chaque résultat explicités.
Les preuves locales, les résultats fonctionnels et les énoncés
conditionnels maintiennent ici leur portée propre.
